% ======================================================================
% sec-related.tex -- related work.
% ======================================================================
\section{Related Work}
\label{sec:related}

\noindent\textbf{Algebraic provenance.}
Semiring provenance~\cite{green2007provenance} and its extensions---%
probabilistic~\cite{suciu2011probabilistic},
factorized~\cite{olteanu2015factorized}, and
aggregate~\cite{senellart2018provenance}---operate after the semantic model
is fixed: a single instance with no unresolved ambiguity, i.e.\ support depth
$0$ here. Why-not provenance~\cite{buneman2001why}, causal
responsibility~\cite{meliou2010causality}, and counterfactual
explanations~\cite{halpern2001causes} likewise condition on one outcome.
Determination provenance recovers all of these as the \emph{result-extensional}
special case (Section~\ref{sec:result-extensional}, Example~\ref{ex:gkt}) and
generalizes them by making the commitments that precede derivation explicit
and composable. The responsibility value
(Section~\ref{sec:responsibility}) and its Shapley
refinement~\cite{livshits2021shapley,shapley1953value} are sequential to, not
competing with, tuple-level responsibility~\cite{meliou2010causality}: there
the players are data tuples (attributing the derivation under a resolution);
here they are semantic commitments (attributing the resolution itself). This
is the key distinction from Shapley-in-databases
work~\cite{livshits2021shapley}, which fixes a single instance.

\smallskip
\noindent\textbf{Possible worlds and probabilistic databases.}
Probabilistic databases~\cite{suciu2011probabilistic} place flat uncertainty
over tuple-existence events with no non-commuting layers; lineage is a Boolean
formula over those events. A tuple-independent PDB is exactly the depth-$1$
case (Proposition~\ref{prop:pdb}). The filtration generalizes this to layered,
non-commuting commitments; the PDB tractability
dichotomy~\cite{dalvi2012probabilistic} is an orthogonal axis that applies
\emph{within} each layer (Section~\ref{sec:pdb}). The separating feature is
the filtration itself: two specifications can induce the same world-set and the
same supports yet different filtrations, and only the filtration certifies
that an answer is settled after examining a bounded prefix.

\smallskip
\noindent\textbf{Provenance beyond the monotone fragment.}
Dannert et al.~\cite{dannert2021semiring} extend semiring provenance to
stratified negation within a single model via monus; for stratified programs
our within-determination provenance coincides with theirs, and
Theorem~\ref{thm:monus-elimination} shows sealing reproduces their
\emph{support} behavior under positivity.
Deutch et al.~\cite{deutch2014circuits} and
K\"ohler et al.~\cite{kohler2012declarative} likewise assume a resolved
semantics. Most closely related, Bogaerts et
al.~\cite{bogaerts2026whynot} define why- and why-not-provenance for
$\mbox{Datalog}^\neg$ under both well-founded and stable semantics using
justification theory, and show the associated decision problem is NP-complete.
Their explanations are \emph{relative to a fixed model}: a justification graph
explains why a chosen literal holds or fails in a given model. The two works
measure complementary axes of cost and are most informative together. Bogaerts
et al.\ fix a determination (one model) and give the \emph{computational} cost
of provenance within it, with a uniform justification-graph representation
across the two semantics. Determination provenance holds within-model
provenance fixed as classical positive $K$-provenance and measures the
\emph{determination} axis: the support and its depth---how many layers of
irrevocable commitment settle an atom, and which commitments flip it across
models (Theorem~\ref{thm:negation-general}). The axes are independent: an atom
can be hard to justify inside a model yet robust across all of them (support
depth $0$), or trivial within every model yet contingent across them. We take
the admissible determinations $\Dets$ as given, isolating the determination
axis; their model-relative why-not is, in our terms, a why-not certificate at a
determination, whereas ours is read from the why-support before any model is
fixed. The two also compose: conditioned on a determination the derivation of
each negative literal is reconstructible (the Gelfond--Lifschitz reduct is
positive once the cycle atoms are committed), so a justification can serve as
the within-model annotation $P_D(a)$. Our contribution is the cross-model,
unstratified setting, where multiple stable models coexist, monus is undefined,
and the determination semiring tracks which models support each atom. The
enriched outcome domain
($\bot \OrdO \mathbf{u} \OrdO \{\mathbf{t},\mathbf{f}\}$) is Belnap's knowledge
order~\cite{belnap1977four}; we use only its refinement structure.

\smallskip
\noindent\textbf{Consistent query answering.}
Arenas et al.~\cite{arenas1999consistent} and Wijsen~\cite{wijsen2019certain}
define certain answers as tuples holding in every repair---analogous to
robustness (full support), but over flat repairs rather than layered
determinations. The analogy is algorithmic as well as conceptual: robustness
is a validity question over the support formula, and the bounded-treewidth
tractability of Theorem~\ref{thm:treewidth} mirrors the structural tractability
results in CQA.

\smallskip
\noindent\textbf{Provenance in systems; related frameworks.}
Database provenance systems (Perm~\cite{glavic2008perm},
GProM~\cite{arab2018gpromprovenance}, Smoke~\cite{psallidas2018smoke}) trace
derivations under a fixed execution. Vandevoort et
al.~\cite{vandevoort2025mixed} certify which transactions may run at weaker
isolation---a cross-specification support comparison in our terms. Ameloot et
al.~\cite{ameloot2013relational} characterize coordination-free evaluation;
Hellerstein~\cite{hellerstein2025complexity} develops a complexity theory for
the cost of resolution. We take determinations as given and explain
algebraically how query results depend on them.
